% =============================================================================
% 论文版式集中配置（对照 ../毕业论文格式规范-计算机与信息技术学院.md）
% 正文尽量只写内容与结构，字体字号与页眉页脚等均在本文声明。
% =============================================================================

% ----- 论文元数据（修改封面与个人信息时只改这里） -----
\newcommand{\thesistitleline}{本科生毕业（设计）论文}
\newcommand{\thesissubtitle}{基于自然语言的知识图谱查询与可视化系统研究与实现}
\newcommand{\thesisauthor}{欧阳闻奕}
\newcommand{\thesisstudentid}{2022102069}
\newcommand{\thesiscollege}{计算机与信息技术学院}
\newcommand{\thesismajorline}{软件工程/2026 届}
\newcommand{\thesisadvisor}{丁蕊}
\newcommand{\thesisadvisorrole}{副教授}
\newcommand{\thesisenterpriseadvisor}{}
\newcommand{\thesisdateyear}{2026}
\newcommand{\thesisdatemonth}{04}
\newcommand{\thesisdateday}{25}
\newcommand{\thesislogopath}{img/logo.png}
\newcommand{\thesislogowidth}{15cm}

% ----- 宏包 -----
\usepackage{ctex}
\usepackage{amsmath,amsfonts,amssymb}
\usepackage{graphicx}
\usepackage[subfigure]{tocloft}
\usepackage{svg}
\usepackage{subfigure}
\usepackage{float}
\usepackage[export]{adjustbox}
\usepackage{listings}
\usepackage{xcolor}
\usepackage{bibentry}
\usepackage[numbers,sort&compress]{natbib}
\usepackage{abstract}
\usepackage{url}
\usepackage{bm}
\usepackage{multirow}
\usepackage{booktabs}
\usepackage{longtable}
\usepackage{supertabular}
\usepackage{changepage}
\usepackage{enumerate}
\usepackage{caption}
\usepackage{indentfirst}
\usepackage{chngcntr}
\usepackage{titlesec}
\usepackage{fancyhdr}
\usepackage{times}

% 页边距：与 Word 模板主节 A4 设置对齐（上/下 2.54cm，左/右 3.18cm）
\usepackage[left=3.18cm,right=3.18cm,top=2.54cm,bottom=2.54cm,headheight=14pt]{geometry}

\renewcommand{\baselinestretch}{1.5}

% ----- 全局正文字体：中文宋体小四（西文 Times 由 times 包） -----
\AtBeginDocument{\songti\zihao{-4}}

% ----- 摘要环境（中文摘要正文宋体小四；标题黑体小二号居中） -----
\setlength{\absleftindent}{0pt}
\setlength{\absrightindent}{0pt}
\renewcommand{\abstractname}{摘\quad 要}
\renewcommand{\abstractnamefont}{\centering\heiti\zihao{-2}}
\renewcommand{\abstracttextfont}{\normalfont\songti\zihao{-4}}

% ----- 关键词（摘要后空两行再写；关键词标签黑体，内容宋体小四） -----
\newenvironment{thesiskeywords}{%
  \par\vspace{2\baselineskip}%
  \noindent{\heiti\zihao{-4} 关键词：}\songti\zihao{-4}%
}{\par}

% ----- 英文摘要 -----
\newenvironment{thesisabstracten}{%
  \begin{center}%
  {\bfseries\zihao{-4} Abstract}%
  \end{center}%
  \vspace{\baselineskip}%
  \begingroup
  \zihao{-4}%
  \setlength{\parindent}{2em}%
}{\par\endgroup}

\newcommand{\thesiskeywordsen}[1]{%
  \par\vspace{0.5\baselineskip}%
  \noindent\textbf{Keywords:} #1\par
}

% ----- 章节编号：第 1 章 … -----
\renewcommand{\thesection}{第\arabic{section}章}
\titleformat{\section}[block]{\centering\heiti\zihao{-2}}{\thesection}{0.5em}{}
\titleformat{\subsection}[hang]{\heiti\zihao{-3}}{\thesubsection}{0.5em}{}
\titleformat{\subsubsection}[hang]{\heiti\zihao{4}}{\thesubsubsection}{0.5em}{}
\titlespacing*{\section}{0pt}{1.5ex plus 0.3ex minus 0.2ex}{0.8ex plus 0.2ex}
\titlespacing*{\subsection}{2em}{1.2ex plus 0.3ex minus 0.2ex}{0.6ex plus 0.2ex}
\titlespacing*{\subsubsection}{2em}{1ex plus 0.3ex minus 0.2ex}{0.5ex plus 0.2ex}

% ----- 图、表按章编号：图 1-1、表 1-1；题注黑体五号 -----
\counterwithin{figure}{section}
\counterwithin{table}{section}
\renewcommand{\thefigure}{\arabic{section}-\arabic{figure}}
\renewcommand{\thetable}{\arabic{section}-\arabic{table}}
\renewcommand{\figurename}{图}
\renewcommand{\tablename}{表}
% 图题、表题：黑体五号（10.5pt）
\DeclareCaptionFont{thesiscap}{\heiti\fontsize{10.5pt}{15.75pt}\selectfont}
\captionsetup{
  font=thesiscap,
  labelfont=thesiscap,
  textfont=thesiscap,
  labelsep=space,
}

% ----- 目录标题 -----
\renewcommand{\contentsname}{目录}
% 目录标题：黑体小二号居中（tocloft 双侧 \hfill 夹住标题）
\renewcommand{\cfttoctitlefont}{\hfill\heiti\zihao{-2}}
\renewcommand{\cftaftertoctitle}{\hfill\null\par\vspace{1\baselineskip}}

% ----- 代码 listings -----
\lstdefinelanguage{cypher}{
  morekeywords={MATCH,RETURN,WHERE,AS,WITH,ORDER,LIMIT,OPTIONAL,CREATE,DELETE,MERGE,DETACH,SET,REMOVE,DROP,CALL,UNWIND,UNION,AND,OR,NOT,IN,CONTAINS,STARTS,ENDS,IS,NULL,CASE,WHEN,THEN,ELSE,END,DISTINCT,SKIP},
  sensitive=false,
  morecomment=[l]{//},
  morestring=[b]',
}
\lstdefinelanguage{json}{
  morekeywords={true,false,null},
  sensitive=false,
  morestring=[b]",
  morecomment=[l]{//},
}
\lstset{
  basicstyle=\small\ttfamily,
  breaklines=true,
  frame=single,
  backgroundcolor=\color{gray!10},
  rulecolor=\color{gray!30},
  numbers=none,
  showstringspaces=false,
  tabsize=2,
  aboveskip=10pt,
  belowskip=10pt
}

\setlength{\cftbeforesecskip}{3pt}
\setlength{\cftbeforesubsecskip}{2pt}
\setlength{\cftbeforesubsubsecskip}{1pt}

% ----- 页眉页脚 -----
\newcommand{\thesisheaderline}{牡丹江师范学院本科生毕业论文（设计）}
\fancypagestyle{thesisempty}{%
  \fancyhf{}%
  \renewcommand{\headrulewidth}{0pt}%
}
\fancypagestyle{thesisfront}{%
  \fancyhf{}%
  \fancyhead[C]{\kaishu\zihao{5}\thesisheaderline}%
  \renewcommand{\headrulewidth}{0.4pt}%
}
\fancypagestyle{thesismain}{%
  \fancyhf{}%
  \fancyhead[C]{\kaishu\zihao{5}\thesisheaderline}%
  \fancyfoot[C]{\thepage}%
  \renewcommand{\headrulewidth}{0.4pt}%
}

% ----- 封面 -----
\newcommand{\thesiscover}{%
  \thispagestyle{thesisempty}%
  \begin{center}
    \includegraphics[width=\thesislogowidth]{\thesislogopath}%
    \vspace{0.5\baselineskip}

    {\heiti\zihao{-2}\thesistitleline\par}%
    \vspace{2cm}%

    {\heiti\zihao{-2} 题目：\thesissubtitle\par}%
    \vspace{3cm}%
  \end{center}%
  \begin{flushleft}
    \heiti\zihao{3}%
    \begin{tabbing}
      姓\hspace{1.5em}名：\hspace{1.5em} \= \thesisauthor \quad \\[0.5cm]
      学\hspace{1.5em}号：\> \thesisstudentid \\[0.5cm]
      学\hspace{1.5em}院：\> \thesiscollege \\[0.5cm]
      专业/届别：\> \thesismajorline \\[0.5cm]
      指导教师：\> \thesisadvisor \\[0.5cm]
      职\hspace{1.5em}称：\> \thesisadvisorrole \\[0.5cm]
      \ifx\thesisenterpriseadvisor\empty\else
      行企业指导教师：\> \thesisenterpriseadvisor \\[0.5cm]
      \fi
    \end{tabbing}%
    \vfill
    \begin{center}
      {\heiti\zihao{3}\thesisdateyear 年\thesisdatemonth 月\thesisdateday 日\par}%
    \end{center}%
  \end{flushleft}%
  \clearpage
}

% 后记级标题（结论 / 致谢 / 参考文献）：黑体小二号居中；不用 \section* 以免与 titlesec 冲突
\newcommand{\thesisbackmattertitle}[1]{%
  \begin{center}{\heiti\zihao{-2} #1\par}\end{center}%
  \vspace{\baselineskip}%
}

\newenvironment{thesisconclusion}{%
  \clearpage
  \phantomsection
  \thesisbackmattertitle{结论}%
  \addcontentsline{toc}{section}{结论}%
  \begingroup
  \songti\zihao{-4}%
}{%
  \par\endgroup
}

\newenvironment{thesisacknowledgement}{%
  \clearpage
  \phantomsection
  \thesisbackmattertitle{致谢}%
  \addcontentsline{toc}{section}{致谢}%
  \begingroup
  \songti\zihao{-4}%
}{%
  \par\endgroup
}

\newcommand{\thesisbibliographyheading}{%
  \clearpage
  \phantomsection
  \thesisbackmattertitle{参考文献}%
  \addcontentsline{toc}{section}{参考文献}%
}

% ----- 颜色（按需保留） -----
\newcommand{\red}[1]{\textcolor[rgb]{1.00,0.00,0.00}{#1}}
\newcommand{\blue}[1]{\textcolor[rgb]{0.00,0.00,1.00}{#1}}
\newcommand{\green}[1]{\textcolor[rgb]{0.00,1.00,0.00}{#1}}
\newcommand{\darkblue}[1]{\textcolor[rgb]{0.00,0.00,0.50}{#1}}
\newcommand{\darkgreen}[1]{\textcolor[rgb]{0.00,0.37,0.00}{#1}}
\newcommand{\darkred}[1]{\textcolor[rgb]{0.60,0.00,0.00}{#1}}
\newcommand{\brown}[1]{\textcolor[rgb]{0.50,0.30,0.00}{#1}}
\newcommand{\purple}[1]{\textcolor[rgb]{0.50,0.00,0.50}{#1}}

\usepackage{hyperref}
\hypersetup{colorlinks=true,linkcolor=black}
